\documentclass{book}

\begin{document}

  \section{Lógica Proposicional}
  
  En lógica proposicional existen 2 tipos de proposiciones, las \textbf{atómicas} y las
  \textbf{moleculares}, y decimos que una proposición es verdadera $V$ bajo una valuación $v$
  si al aplicar las reglas semánticas de los conectores, el resultado es $V$.
  
  Para el caso de una proposición atómica $p$ tenemos únicamente 2 posibles valores $V$ y $F$.
  Para el caso de las moleculares cambia para cada conector que indica cómo interactúan.
  El valor de verdad está definido en la semántica que propaga. Los casos en que una valuación
  es $F$ se conocen como \textbf{caso de exclusión}.
  
  \note[title={Valuación}]{La valuación $V$ es determinada por las reglas del modelo al que se aplica.}
  
  \subsection{Negación}
  
  \begin{minipage}{0.55\textwidth}
  \begin{definition}[title={Negación $\sim p$}]
  La molécula más simple, puesto que niega el valor de valuación de $p$.
  \end{definition}
  \end{minipage}
  \begin{minipage}{0.35\textwidth}
  \begin{tabular}{c|c}
  $p$ & $\sim p$ \\
  \hline
  V & F \\
  F & V \\
  \end{tabular}
  \end{minipage}
  
  \subsection{Conjunción}
  
  \begin{minipage}{0.55\textwidth}
  \begin{definition}[title={Conjunción $p \land q$}]
  Este conector solo permite el par ordenado $(V, V)$ y excluye al resto.
  Decimos entonces que $p \land q$ se puede representar por pares ordenados
  en un conjunto sin las exclusiones, ie:
  \begin{equation}
  p \land q = \{(V, V)\}
  \end{equation}
  \end{definition}
  \end{minipage}
  \begin{minipage}{0.35\textwidth}
  \begin{tabular}{c|c|c}
  $p$ & $q$ & $p \land q$ \\
  \hline
  V & V & V \\
  V & F & F \\
  F & V & F \\
  F & F & F \\
  \end{tabular}
  \end{minipage}
  
  \subsection{Disyunción}
  
  \begin{minipage}{0.55\textwidth}
  \begin{definition}[title={Disyunción $p \lor q$}]
  Para este conector únicamente se excluye $(F, F)$. Por lo tanto se puede representar que:
  \begin{equation}
  p \lor q = \{(V,V),\ (V,F),\ (F,V)\}
  \end{equation}
  \end{definition}
  \end{minipage}
  \begin{minipage}{0.35\textwidth}
  \begin{tabular}{c|c|c}
  $p$ & $q$ & $p \lor q$ \\
  \hline
  V & V & V \\
  V & F & V \\
  F & V & V \\
  F & F & F \\
  \end{tabular}
  \end{minipage}
  
  En lógica proposicional no se busca describir el mundo físico, ie, la lógica no entiende
  causalidad ni temporalidad, sino que son un conjunto de reglas.
  
  \subsection{Implicación}
  
  \begin{definition}[title={Implicación $p \to q$}]
  Este conector no justifica que si ocurre $p$ causa $q$, ni que para que ocurra $q$ deba
  ocurrir $p$. No se evalúa de forma secuencial, sino que cada uno tiene su valor.
  La única exclusión que tiene es el caso $(V, F)$.
  \end{definition}
  
  \begin{minipage}{0.4\textwidth}
  \begin{tabular}{c|c|c}
  $p$ & $q$ & $p \to q$ \\
  \hline
  V & V & V \\
  V & F & F \\
  F & V & V \\
  F & F & V \\
  \end{tabular}
  \end{minipage}
  \begin{minipage}{0.55\textwidth}
  Dicho de otra forma, si $p$ es falso, significa que la hipótesis no se cumple,
  por tanto no se considera el valor de $q$. La regla de exclusión dice:
  
  \alert{No existe un caso en que $p$ sea $V$ y $q$ sea $F$}
  
  \alert{Si no hay promesa, no hay incumplimiento}
  \end{minipage}
  
  El caso $(F, F)$ puede parecer contraintuitivo; ambas falsas debería ser falsa la proposición
  en general, pero esto no es así porque no se busca reflejar o explicar una situación del mundo físico.
  
  \subsection{Equivalencia lógica}
  
  Antes de continuar con otros elementos, debemos aclarar la idea de la \textbf{equivalencia
  lógica}, ya que en ocasiones resulta más sencillo abordar una equivalencia a toda la
  proposición original.
  
  \begin{definition}[title={Equivalencia}]
  Decimos que dos proposiciones son equivalentes si tienen los mismos casos de exclusión.
  \end{definition}
  
  De acuerdo con esto podemos hallar una equivalencia para $p \to q$. Un caso de exclusión
  es $(V, F)$, ie, $p = V$ y $q = F$, lo cual se expresa como $p \land \sim q$. Sin embargo,
  esto es el caso que \textbf{no} está permitido, por lo tanto negamos esa expresión y
  hallamos la expresión adecuada:
  
  \begin{tabular}{c|c|c|c|c|c}
  $p$ & $q$ & $p \to q$ & $\sim q$ & $p \land \sim q$ & $\sim(p \land \sim q)$ \\
  \hline
  V & V & V & F & F & V \\
  V & F & F & V & V & F \\
  F & V & V & F & F & V \\
  F & F & V & V & F & V \\
  \end{tabular}
  
  \alert{De modo que los valores de exclusión son aquellos en los que el conector es falso
  por definición semántica.}
  
  \subsection{Bicondicional}
  
  \begin{definition}[title={Bicondicional $p \leftrightarrow q$}]
  Es verdad cuando $p$ y $q$ tienen el mismo valor de verdad en todas las interpretaciones.
  Este conector es una síntesis útil para una composición de implicaciones:
  \begin{equation}
  p \leftrightarrow q \equiv (p \to q) \land (q \to p)
  \end{equation}
  De forma que es verdadera cuando coinciden $p$ y $q$:
  \begin{equation}
  p \leftrightarrow q \equiv (p \land q) \lor (\sim p \land \sim q)
  \end{equation}
  \end{definition}
  
  \begin{tabular}{c|c|c|c|c}
  $p$ & $q$ & $p \to q$ & $q \to p$ & $(p \to q) \land (q \to p)$ \\
  \hline
  V & V & V & V & V \\
  F & V & V & F & F \\
  V & F & F & V & F \\
  F & F & V & V & V \\
  \end{tabular}
  
  \subsection{Demostración}
  
  Las equivalencias las escribimos por análisis, es decir, sabemos cómo se comportan los
  conectores. Recordemos que la lógica proposicional no busca explicar el mundo físico;
  para probar que una proposición es válida, las dos formas más comunes son
  \textbf{prueba directa} o \textbf{por contraejemplo}.
  
  \begin{block}[title={Estrategia general}]
  Para probar una proposición, se debe probar que el caso donde se define como $F$
  no puede suceder, es decir, mostrar que los casos de exclusión no son válidos.
  \end{block}
  
  \begin{proposition}[title={Si $n$ es par, entonces $n^2$ también es par}]
  \end{proposition}
  
  \begin{proof}[title={}]
  Sea $p$: "$n$ es par" y $q$: "$n^2$ es par". Se tiene $p \to q$.
  Se prueba que $p \land \sim q$ no puede suceder, ie, que si $n$ es par, $n^2$ no puede ser impar.
  
  Sea $n = 2k$:
  \begin{align}
  n^2 &= (2k)^2 \\
      &= 4k^2 \\
      &= 2(2k)
  \end{align}
  
  Por lo tanto $n^2$ es par, lo cual contradice $\sim q$. De esta forma, $p \land \sim q$
  no puede ocurrir, luego $p \to q$ es una tautología. \qed
  \end{proof}
  
  \begin{proposition}[title={Si el triángulo $ABC$ tiene 2 ángulos de $60°$ entonces es equilátero}]
  \end{proposition}
  
  \begin{minipage}{0.5\textwidth}
  \begin{proof}[title={Construcción directa}]
  En un espacio euclídeo, todo triángulo tiene $180°$, de modo que el tercer ángulo es $60°$ pues:
  \begin{align}
  x + 2(60) &= 180 \\
  x &= 60°
  \end{align}
  Luego, $ABC$ es equilátero. \qed
  \end{proof}
  \end{minipage}
  \begin{minipage}{0.45\textwidth}
  Si bien no se demuestra directamente la exclusión, se recurre a ella pues:
  \begin{align}
  &\text{ABC tiene 2 ángulos de } 60° \to \text{ABC es equilátero} \\
  &p \to q
  \end{align}
  \end{minipage}
  
  Para este caso no se prueba $p \land \sim q$ directamente, sino que tomando $p$ como
  hipótesis, el modelo (geometría euclídea) predice que $\sim q$ es imposible, de modo que
  $p \land \sim q$ no puede ocurrir.
  
  Por ello, para probar la validez del bicondicional se suele probar $p \to q$ y posterior
  $q \to p$, pues los valores excluidos son $(V,F)$ y $(F,V)$ para cada implicación, de modo
  que $p \leftrightarrow q$ excluye ambos. Formalmente:
  
  \begin{equation}
  p \leftrightarrow q \equiv \sim((\sim p \land q) \lor (p \land \sim q))
  \end{equation}
  
  \note[title={Patrón tautológico}]{Nótese que se está aplicando la tautología como patrón,
  de forma que puede sustituir o simplificar proposiciones previamente probadas.
  Este mecanismo fue usado por Euclides en sus libros clásicos.}
  
  \subsection{Contraposición}
  
  \begin{definition}[title={Contraposición}]
  Si se asume que $q$ es $F$, entonces $p$ también debe serlo, puesto que $p \to q$
  excluye el caso $p = V$ y $q = F$:
  \begin{equation}
  p \to q \equiv \sim q \to \sim p
  \end{equation}
  \end{definition}
  
  \begin{tabular}{c|c|c|c|c|c}
  $p$ & $q$ & $p \to q$ & $\sim p$ & $\sim q$ & $\sim q \to \sim p$ \\
  \hline
  V & V & V & F & F & V \\
  V & F & F & F & V & F \\
  F & V & V & V & F & V \\
  F & F & V & V & V & V \\
  \end{tabular}
  
  \subsection{Leyes de De Morgan}
  
  \begin{theorem}[title={Leyes de De Morgan}]
  \begin{align}
  \sim(p \land q) &\equiv \sim p \lor \sim q \\
  \sim(p \lor q)  &\equiv \sim p \land \sim q
  \end{align}
  \end{theorem}
  
  Para que no ocurra $p \land q$, al menos una debe ser $F$.
  Ambas deben ser $F$ para que $p \lor q$ no se cumpla.
  
  \subsection{Formas Normales}
  
  Algo que puede ayudarnos a verificar factibilidad es convertir una proposición en su
  \textbf{forma normal}; esta es una forma estandarizada de escribir proposiciones usando
  solo ciertos conectores. Toda proposición puede representarse por una FN.
  
  \begin{minipage}{0.45\textwidth}
  \begin{definition}[title={FND — Forma Normal Disyuntiva}]
  Es una conjunción de disyunciones:
  \begin{equation}
  (p_1 \lor q_1) \land (p_2 \lor q_2) \land \cdots
  \end{equation}
  \end{definition}
  \end{minipage}
  \begin{minipage}{0.45\textwidth}
  \begin{definition}[title={FNC — Forma Normal Conjuntiva}]
  Es una disyunción de conjunciones:
  \begin{equation}
  (p_1 \land q_1) \lor (p_2 \land q_2) \lor \cdots
  \end{equation}
  \end{definition}
  \end{minipage}
  
  Para convertir una proposición a una forma normal, se siguen estos 3 pasos:
  
  \begin{enumerate}
  \item Se eliminan las implicaciones.
  \item Se expanden las negaciones usando las leyes de De Morgan, con la finalidad de
  poder agrupar linealmente.
  \item Se distribuyen los conectores.
  \end{enumerate}
  
  Por ejemplo, para $p \to (q \land r)$:
  \begin{align}
  p \to (q \land r) &\equiv \sim p \lor (q \land r) \\
                    &\equiv (\sim p \lor q) \land (\sim p \lor r)
  \end{align}
  
  \subsection{Reglas de Inferencia}
  
  Podemos simplificar la notación asignando nombres a ciertas proposiciones; estos son
  conocidos como \textbf{reglas de inferencia}.
  
  \begin{minipage}{0.45\textwidth}
  \begin{definition}[title={Modus Ponens}]
  La restricción se propaga:
  \begin{align}
  &p \to q \nonumber\\
  &p \nonumber\\
  &\therefore\ q \nonumber
  \end{align}
  \end{definition}
  \end{minipage}
  \begin{minipage}{0.45\textwidth}
  \begin{definition}[title={Modus Tollens}]
  La restricción es $p \land \sim q$ y $\sim q = V$, no puede ocurrir $p = V$,
  solo queda $\sim p$:
  \begin{align}
  &p \to q \nonumber\\
  &\sim q \nonumber\\
  &\therefore\ \sim p \nonumber
  \end{align}
  \end{definition}
  \end{minipage}
  
  \begin{minipage}{0.45\textwidth}
  \begin{definition}[title={Silogismo Hipotético}]
  Una composición de restricciones:
  \begin{align}
  &p \to q \nonumber\\
  &q \to r \nonumber\\
  &\therefore\ p \to r \nonumber
  \end{align}
  \end{definition}
  \end{minipage}
  \begin{minipage}{0.45\textwidth}
  \begin{definition}[title={Silogismo Disyuntivo}]
  Un \textit{o} excluyente:
  \begin{align}
  &p \lor q \nonumber\\
  &\sim p \nonumber\\
  &\therefore\ q \nonumber
  \end{align}
  \end{definition}
  \end{minipage}
  
  \subsection{Argumento Válido}
  
  \begin{definition}[title={Argumento}]
  Un argumento tiene la forma $(P_1, P_2, \ldots) \therefore q$, o formalmente
  $(P_1 \land P_2 \land \cdots) \to q$. Diremos que es \textbf{válido} si no existe
  una valuación donde las premisas son verdaderas y la conclusión es falsa.
  \end{definition}
  
  Toda regla de inferencia se puede justificar mostrando que determinada fórmula es una
  tautología. Para verificar el modus ponens, construimos $((p \to q) \land p) \to q$:
  
  \begin{tabular}{c|c|c|c|c}
  $p$ & $q$ & $p \to q$ & $(p \to q) \land p$ & $((p \to q) \land p) \to q$ \\
  \hline
  V & V & V & V & V \\
  V & F & F & F & V \\
  F & V & V & F & V \\
  F & F & V & F & V \\
  \end{tabular}
  
  De este modo, si tenemos $(((p \to q) \land (q \to r)) \land p) \to r$, podemos reemplazar
  todo el primer bloque por $r$ sin perder verdad.
  
  \begin{definition}[title={Verdad bajo valuación}]
  Una fórmula es verdad bajo una valuación $v$ si al evaluarla con las reglas semánticas
  de los conectores el resultado es $V$. Un argumento es válido si no existe una valuación
  donde todas las premisas sean verdaderas y las conclusiones sean falsas.
  \end{definition}

\end{document}

